\chapterimage{figs/chapterhead.png}
\chapter{First steps in the graphical interface}
\label{chap:primeiros_passos_interface_grafica}

\section{Introduction}
\label{sec:cap3_introducao}

Once the \textit{GenESyS} graphical application is running, the user actually enters the simulator. It is from that moment on that the system stops being just a compiled project and becomes a work environment. This chapter has precisely this function: transforming the first execution of the GUI into a guided first contact with the interface.

The proposal here is not yet to model in detail or run complete examples. Before that, the reader needs to recognize the logic of the main window, understand which areas of the interface are most important, understand where the essential commands are and build a sufficiently clear notion of the workspace that \textit{GenESyS} offers. In other words, this chapter teaches the user how to locate themselves within the application.

This step is more important than it may seem. A rich interface, with multiple menus, tabs, trees, panels and simulation commands, can generate the impression of excessive complexity when the user tries to understand it at once. The best strategy is different: first recognize the global structure of the application, then identify its functional groups and, only then, start working with concrete models. It is exactly this route that will be followed here.

\section{The Main Window as a Workspace}
\label{sec:cap3_janela_principal_ambiente}

\textit{GenESyS}'s graphical interface is organized around a main window that centralizes modeling, inspection, simulation, debugging, and reporting operations. This means that the user is not faced with an application with a single panel or a minimal set of isolated commands. Instead, the GUI was designed as a complete work environment in which different tasks can be performed through different views of the same model.

From a usage point of view, this main window should be understood as the space where five fundamental dimensions of working with the simulator are articulated:

\begin{enumerate}
\item model management;
\item graphic editing;
\item parameterization and inspection;
\item running the simulation;
\item observation of results and system behavior.
\end{enumerate}

This organization is especially important because the user does not work with the model just before execution or just after it. The model is continually built, adjusted, inspected, and reevaluated. The GUI was structured to support exactly this cycle.

\section{Opening the Application and First Visual Reading}
\label{sec:cap3_abertura_e_leitura_visual}

When starting the graphical application, the user must resist the temptation to immediately click on all available menus and commands. The best way is a guided visual reading of the main window. This reading can be done in very simple steps.

First, look at the menu bar and notice that it organizes work into functional categories. In general, model actions, editing, visualization, simulation, tools and auxiliary information are present. Second, notice the central area of ​​the interface, which is the main graphical representation space of the model. Third, note the existence of tabs, side or bottom panels and navigation trees, which complement the central view. Fourth, note that the application was not designed just to edit a diagram, but to monitor the entire life of the model, including its simulation and inspection.

This first reading is enough to draw an important conclusion: the \textit{GenESyS} GUI is not a simple graphical editor. It is an integrated work interface.

\section{Interface Functional Groups}
\label{sec:cap3_grupos_funcionais_interface}

The most productive way to understand the GUI is to divide it into functional groups. Instead of memorizing isolated commands, the user should ask: what kind of work does each part of the interface organize?

\subsection{Template Commands}
\label{subsec:cap3_comandos_modelo}

The first functional group is model commands. They are what allow you to start a new job, open an already saved model, save the current model, close it, check its consistency and obtain information about it. These commands structure the most basic cycle of using the tool.

In practice, this means that the user must recognize at the outset where the actions responsible for:
\begin{itemize}
\item create a new model;
\item open an existing model;
\item save work;
\item close the current model;
\item check the model before simulating.
\end{itemize}

These commands will be used repeatedly throughout the first part of the book.

\subsection{Editing Commands}
\label{subsec:cap3_comandos_edicao}

The second group is editing commands. This includes actions such as undo, redo, copy, paste, delete, group, ungroup and, in some cases, organization and alignment commands. This set reveals that the interface was designed for effective editing of the model, and not just for passive observation.

Even if the user is not yet building new models, it is important to realize right away that the GUI supports modification of graphic and structural elements. This will be useful when the example models start to change, even slightly.

\subsection{View Commands}
\label{subsec:cap3_comandos_visualizacao}

Another important group is view commands. They include adjustments such as zoom, grid, guides and other features that facilitate the visual organization of the graphic space. On first contact, the user does not need to explore all these resources in detail, but must recognize them as tools to support reading and editing the model.

In particular, zooming is often useful right from the start, because different models may require different observation scales.

\subsection{Simulation Commands}
\label{subsec:cap3_comandos_simulacao}

A central group of the GUI is made up of commands linked to running the simulation. These typically include starting, pausing, resuming, stopping, step-by-step, and configuring the simulation. These commands show that the application not only organizes the model as a structure, but also monitors its execution.

In this chapter, the reader must only recognize the existence and location of these commands. Its effective use will be deepened when we start working with concrete models. Still, it's important to stick to the idea that the same GUI used to open and edit the model will also be used to conduct your simulation.

\subsection{Inspection and Support Commands}
\label{subsec:cap3_comandos_inspecao_apoio}

In addition to the previous groups, the interface also offers inspection and support mechanisms: property panels, component trees and \textit{data definitions}, debug areas, report tabs and other auxiliary views. The presence of these resources reinforces an essential point: \textit{GenESyS} was not designed just to design models, but to study them in execution.

\begin{table}[htbp]
\centering
\caption{Main Functional Groups of the GenESyS GUI.}
 \label{tab:cap3_grupos_funcionais}
\small
 \begin{tabular}{|p{4.0cm}|p{7.4cm}|}
\hline
\textbf{Functional Group} & \textbf{Role in the User's Work} \\
\hline
Template & Create, open, save, close, check and organize work with templates\\
\hline
Editing & Changing and rearranging model elements \\
\hline
View & Adjust scale, guides and visual organization of the graphics area \\
\hline
Simulation & Start, pause, resume, stop and track execution \\
\hline
Inspection and support & Examine properties, model elements, results and auxiliary information \\
\hline
 \end{tabular}
\end{table}

\section{The Graphics Area of the Model}
\label{sec:cap3_area_grafica_modelo}

The graphics area is the visual center of working with \textit{GenESyS}. It is where the model appears in a diagrammatic form and it is where most of the user's attention tends to focus when reading and editing the models. At first glance, this area should be understood as the main space where the model's flow becomes visible.

This is important because the user needs to start reading the simulator by the way it shows the model, and not just by the way the program organizes its menus. The graphic area allows you to view:
\begin{itemize}
\item the components present in the model;
\item the connections between these components;
\item the general logic of the modeled flow;
\item the visual distribution of elements within the workspace.
\end{itemize}

Even before opening a concrete model, the user should already perceive this area as the core of the simulator's visual experience.

\subsection{Graphic area is not just drawing}
\label{subsec:cap3_area_grafica_nao_apenas_desenho}

An important point is that the graphics area should not be confused with a simple drawing editor. What is organized there are not just geometric shapes, but the representation of the model that will later be simulated. In other words, graphic space is also semantic space. Each visible element there tends to correspond to a relevant modeling entity.

This observation helps the user to look at the interface in the correct way. Instead of thinking that he is just manipulating visual objects, he must understand that he is interacting with the operational representation of the model.

\section{Trees, Panels, and Helper Editors}
\label{sec:cap3_arvores_paineis_editores}

In addition to the central graphics area, the \textit{GenESyS} GUI offers other elements that are equally important for everyday work. These elements complement the visual reading of the model and make it possible to inspect it in a richer way.

\subsection{Navigation Trees}
\label{subsec:cap3_arvores_navegacao}

Navigation trees help the user locate system and model elements. Depending on the context, they may show:
\begin{itemize}
\item \textit{plugins};
\item model components;
\item \textit{data definitions};
\item other auxiliary structures associated with the current work.
\end{itemize}

At first glance, these trees should be seen as guidance mechanisms. They help answer the question: What elements are available or present in the model I'm working with?

\subsection{Property Editor}
\label{subsec:cap3_editor_propriedades}

The property editor is one of the most important tools in the GUI. It is through this that the user stops just observing the existence of an element and starts dealing with its attributes, parameters and configuration options. In other words, this is where the template starts to become effectively editable.

In practice, this means that the selection of an element in the graphics area or in a navigation tree tends to be associated with the display of corresponding properties. The next chapters, this mechanism will appear continuously, because almost every relevant modification of a model goes through some type of parameterization.

\subsection{Complementary tabs and views}
\label{subsec:cap3_abas_visoes_complementares}

The presence of central tabs and panels organized by tabs shows that the GUI was built to support multiple perspectives on the same model. Some of these perspectives are visual, others are textual, others are focused on execution, debugging, or presenting results. The user does not need to master all of this now, but must retain the idea that the model is not observed only through a single visual surface.

This multiplicity of views will be especially important when, later on, we introduce the \textit{Simulang} tab as a complementary textual representation to the graphical model.

\section{Minimal User Guidance Flow in the GUI}
\label{sec:cap3_fluxo_minimo_orientacao}

In order to familiarize yourself with the GUI to be productive, it is advisable to adopt a short guideline. Instead of randomly exploring the main window, the user can follow this sequence:

\begin{enumerate}
\item find template commands;
\item identify the central graphics area;
\item locate auxiliary trees and panels;
\item find the property editor;
\item recognize basic simulation commands;
\item note the existence of additional tabs.
\end{enumerate}

This script has an important advantage: it prevents the richness of the interface from being perceived as clutter. By following a simple order, the user transforms the main window into something intelligible and progressively explorable.

\section{What Is Not Necessary to Master Now}
\label{sec:cap3_o_que_nao_dominar_agora}

A successful first contact with the GUI does not require mastering the entire application. It is not necessary, for example:
\begin{itemize}
\item use all editing commands;
\item understand all debug panels;
\item interpret all interface trees;
\item know in depth all types of available elements;
\item also explore all of the model's languages ​​and representations.
\end{itemize}

On the contrary, trying to master everything immediately tends to produce noise. The most efficient learning occurs in layers. First, the user understands the general structure of the interface. Then, he starts working with concrete models. Then, it begins to report the GUI with execution, results and textual language. This is the route adopted in this manual.

\section{Good Practices When First Encountering the Interface}
\label{sec:cap3_boas_praticas_primeiro_contato}

There are some simple practices that make the initial experience with the GUI much more beneficial.

The first is to always work with a concrete objective. Instead of ``exploring the interface'', try to locate something specific: the open model command, the graphics area, the properties panel or the simulation commands.

The second is not to confuse quantity of resources with difficulty of use. The \textit{GenESyS} interface is broad because it organizes several phases of working with models. This does not mean that all resources need to be used simultaneously.

The third is to look at the interface as an integrated environment. Menus, graphics area, trees, properties, simulation and reports are not independent parts. They form a single workflow.

The fourth is to accept that some parts of the GUI will only make sense when a model is actually loaded. That's why the next chapter will go directly into the example models.

\section{Final Considerations}
\label{sec:cap3_consideracoes_finais}

This chapter introduced the \textit{GenESyS} graphical interface as the user's main working environment. Instead of treating the GUI as a dispersed set of commands, the text sought to show its organizational logic: the main window brings together resources to manage models, edit them, parameterize them, simulate them and observe them from different perspectives. This understanding is essential because it prevents the user from reducing the application to a simple graphical editor or, at the opposite extreme, perceiving it as an excessively complex and cluttered interface.

With this guidance in place, the next natural step is to start working on concrete models. This is when the graphical interface really takes on operational meaning. Therefore, the following chapter will deal with the example models already available in the project and their initial use in the GUI, allowing the reader to open, inspect, execute and modify real simulation cases.

Test your understanding of this content by answering the following questions:

\begin{enumerate}
\item Why should the \textit{GenESyS} GUI be understood as an integrated work environment, and not just as a graphical editor?

\item What are the most important functional groups of the interface in a first contact?

\item What is the function of the model's graphical area within the simulator usage experience?

\item Why isn't it necessary to immediately master all the application's menus, trees and panels?
\end{enumerate}
